\documentclass[CJK]{beamer}
\usepackage{CJKutf8}
\usepackage{beamerthemesplit}
\usetheme{Malmoe}
\useoutertheme[footline=authortitle]{miniframes}
\usepackage{amsmath}
\usepackage{amssymb}
\usepackage{graphicx}
\usepackage{color}
\usepackage{slashed}
\usepackage{simplewick}
\graphicspath{{../figures/}}
\def\be{\begin{equation}}
\def\ee{\nonumber\end{equation}}
\def\bea{\begin{eqnarray}}
\def\eea{\nonumber\end{eqnarray}}
\def\ii{{\dot{\imath}}}
\def\bch{\begin{CJK}{UTF8}{gbsn}}
\def\ech{\end{CJK}}
\def\bex{\begin{minipage}{0.3\textwidth}\includegraphics[width=1in]{jugelizi.png}\end{minipage}\begin{minipage}{0.6\textwidth}}
\def\eex{\end{minipage}}
\def\chtitle#1{\frametitle{\bch#1\ech}}
\def\skipline{{\vskip0.1in}}
\def\skiplines{{\vskip0.2in}}
\def\lagr{{\mathcal{L}}}
\def\hamil{{\mathcal{H}}}
\def\vecv{{\mathbf{v}}}
\def\vecx{{\mathbf{x}}}
\def\veck{{\mathbf{k}}}
\def\vecp{{\mathbf{p}}}
\def\vecn{{\mathbf{n}}}
\def\vecA{{\mathbf{A}}}
\def\vecP{{\mathbf{P}}}
\def\vecsigma{{\mathbf{\sigma}}}
\def\hatJn{{\hat{J_\vecn}}}
\def\hatJx{{\hat{J_x}}}
\def\hatJy{{\hat{J_y}}}
\def\hatJz{{\hat{J_z}}}
\def\hatj#1{\hat{J_{#1}}}
\def\hatphi{{\hat{\phi}}}
\def\hatq{{\hat{q}}}
\def\hatpi{{\hat{\pi}}}
\def\vel{\upsilon}
\def\Dint{{\mathcal{D}}}
\def\adag{{\hat{a}^\dagger}}
\def\bdag{{\hat{b}^\dagger}}
\def\cdag{{\hat{c}^\dagger}}
\def\ddag{{\hat{d}^\dagger}}
\def\hata{{\hat{a}}}
\def\hatb{{\hat{b}}}
\def\hatc{{\hat{c}}}
\def\hatd{{\hat{d}}}
\def\hatN{{\hat{N}}}
\def\hatH{{\hat{H}}}
\def\hatp{{\hat{p}}}
\def\Fup{{F^{\mu\nu}}}
\def\Fdown{{F_{\mu\nu}}}
\def\newl{\nonumber \\}
\def\SIkm{\mathrm{km}}
\def\SIyr{\mathrm{yr}}
\def\SIGyr{\mathrm{Gyr}}
\def\SIeV{\mathrm{eV}}
\def\SIGeV{\mathrm{GeV}}
\def\SIm{\mathrm{m}}
\def\SIcm{\mathrm{cm}}
\def\SIJ{\mathrm{J}}
\def\SIs{\mathrm{s}}
\def\SIkg{\mathrm{kg}}
\def\SIg{\mathrm{g}}
\def\vece{\mathrm{e}}
\def\bmat#1{\left(\begin{array}{#1}}
\def\emat{\end{array}\right)}
\def\bcase#1{\left\{\begin{array}{#1}}
\def\ecase{\end{array}\right.}
\def\calM{{\mathcal{M}}}
\def\calT{{\mathcal{T}}}
\def\calR{{\mathcal{R}}}
\def\barpsi{\bar{\psi}}
\def\baru{\bar{u}}
\def\barv{\bar{\upsilon}}
\def\bmini#1{\begin{minipage}{#1\textwidth}}
\def\emini{\end{minipage}}
\def\qeq{\stackrel{?}{=}}
\def\torder#1{\mathcal{T}\left(#1\right)}
\def\rorder#1{\mathcal{R}\left(#1\right)}


\title{Quantum Field Theory I \\ Lesson 01 - Introduction}
  \author{}
  \date{}


\begin{document}

\begin{frame}
 
\begin{center}
\begin{Large}
\bch
量子场论 I \\
第一课：坐标系和张量场简介

{\vskip 0.3in}

授课: 黄志琦

\ech
\end{Large}
\end{center}

\vskip 0.2in

\bch
教材：《简明量子场论》，王正行，北京大学出版社
\ech

\bch
课件下载
\ech
https://github.com/zqhuang/SYSU\_QFTI

\end{frame}

\begin{frame}
\frametitle{\bch 评分 \ech (Evaluation)}
\bch
总分 = 作业30分 + 课堂表现20分  + 期末考试50分
\ech
\end{frame}

\begin{frame}
\frametitle{\bch 互相简单介绍 \ech}
\bch
自我介绍：
\begin{itemize}
\item{名字}
\item{物理基础：高等数学，线性代数，理论力学，狭义相对论，量子力学这五门课的学习情况（这些是什么鬼?学过大部分忘了?学过大部分还会?精通?）}
\end{itemize}
\ech
\end{frame}

\begin{frame}
  \chtitle{第一章要掌握的内容}
\bch   
\begin{itemize}
\item{自然单位制和量纲}
\item{坐标，度规和张量}
\item{作用量和拉氏密度的概念}
\item{Euler-Lagrange方程的推导}
\item{Noether定理的推导}
\end{itemize}

\ech
\end{frame}

\begin{frame}
\frametitle{\bch 自然单位制和量纲 \ech}
\bch 自然单位制 \ech $c = \hbar  = 1$

\bch
长度量纲 = 时间量纲 = 质量量纲$^{-1}$ = 能量量纲$^{-1}$ 
\ech

{\vskip 0.2in}
\begin{minipage}{0.3\textwidth}
\includegraphics[width=1in]{jugelizi.png}
\end{minipage}
\begin{minipage}{0.6\textwidth}
\bch
力的量纲 \\
= 质量量纲 $\times$ 长度量纲/时间量纲$^2$ \\
= 质量量纲$^2$
\ech
\end{minipage}

\end{frame}

\begin{frame}
\frametitle{\bch 思考题\ech}
\bch
\addfig{0.5}{think.jpg}

下面物理量的量纲是质量的多少次方？

速度$v$

角速度$\omega$

角动量$L$

加速度$a$

能量密度$\rho$

压强$p$

牛顿引力常数$G$

\ech
\end{frame}


\begin{frame}
\frametitle{\bch 思考题 \ech}
\bch
$1\SIkg$的倒数是多少$\SIcm$?

\bcenter
\lfig{1}{blackq.jpg}

根本不知道你在说什么.jpg
\ecenter
\ech
\end{frame}



\begin{frame}
\frametitle{\bch 时空坐标的标记以及一些约定 \ech}
\bch
\begin{itemize}
\item{时空用坐标$(x^0, x^1, x^2, x^3)$来标记，其中$x^0$为时间坐标，有时也写作$t$；省略上标的坐标$x$是$(x^0, x^1, x^2, x^3)$的简写。}
\item{希腊字母指标$\mu, \nu, \rho, \sigma,\ldots$ 默认跑遍0，1，2，3}
\item{拉丁字母指标$i, j,  k, l,\ldots$ 默认跑遍1，2，3}
\item{黑体表示三维空间矢量，例如 $\mathbf{dx} \equiv (dx^1, dx^2, dx^3)$}
\end{itemize}
\ech
\end{frame}


\begin{frame}
\frametitle{\bch 时空坐标的标记以及一些约定 \ech}
\bch
\begin{itemize}
\item{爱因斯坦求和规则：重复出现的指标默认求和。例如：$dx^\mu dx_{\mu}$若无特殊说明等价于$\sum_{\mu = 0}^3 dx^\mu dx_{\mu}$} 
\item{空间坐标的偏微分$\frac{\partial}{\partial x^\mu}$经常简写成$\partial_\mu$. 省略下标的偏微分符号$\partial$是$(\partial_0,\partial_1,\partial_2,\partial_3)$的简写。算符$\nabla$则仅是空间偏微分$(\partial_1, \partial_2, \partial_3)$的简写。}
\item{克罗内克符号(Kronecker delta) $\delta^\mu_{\ \nu}$当$\mu = \nu$时为1，否则为0.}
\end{itemize}
\ech

\end{frame}


\begin{frame}
\chtitle{度规(metric)}
\bch 时空间隔元的定义：

$ds^2 = g_{\mu\nu} dx^\mu dx^\nu$

\ech
\end{frame}


\begin{frame}
\frametitle{\bch 度规 (metric) \ech}
\bch 时空间隔元：\ech
\bch 狭义相对论的平直时空，也叫闵氏空间(Minkowski spacetime)的度规 \ech

\begin{equation} 
g_{\mu\nu} =  \left( \begin{array}{rrrr} 1 & 0 & 0 & 0 \\ 0 & -1 & 0 & 0 \\ 0 & 0 & -1 & 0 \\ 0 & 0 & 0 & -1 \end{array} \right) \nonumber
\end{equation}

\bch 这样时空间隔元可以简写成 $ds^2 = dt^2 - \mathbf{dx}^2$ \ech

\end{frame}


\begin{frame}
\frametitle{\bch 一般性张量(tensor)的定义 \ech}
\bch
一般性张量在坐标变换$x \rightarrow \tilde{x}$下的变化：


$$\tilde{T}^{\mu\nu\ldots}_{\ \ \alpha\beta\ldots} = \left(\frac{\partial \tilde{x}^\mu}{\partial x^\rho}\frac{\partial \tilde{x}^\nu}{\partial x^\sigma} \frac{\partial x^\gamma }{\partial \tilde{x}^\alpha}\frac{\partial x^\lambda}{\partial \tilde{x}^\beta} \ldots \right)T^{\rho\sigma\ldots}_{\ \ \gamma\lambda\ldots} $$

\skiplines

(上个学年很多同学看到这里就退课了\bye)
\ech
\end{frame}

\begin{frame}
\frametitle{\bch 思考题 \ech}
\bch
如果$\phi$是一个标量（0阶张量，在坐标变换下不变），时空度规为$g_{\mu\nu}$，时空度规的逆矩阵为$g^{\mu\nu}$。试证明：
\begin{itemize}
\item{任何常数都是标量}
\item{时空微元$dx^\mu$是张量（一阶张量，也称为矢量）}
\item{$g_{\mu\nu}$本身是张量 （二阶张量）}
\item{$\partial_\mu\phi$是张量（矢量）}
\item{克罗内克符号$\delta^{\mu}_{\ \nu}$是张量}
\item{度规的逆矩阵$g^{\mu\nu}$是张量}
\item{如果定义$dx_{\mu} \equiv g_{\mu\nu} dx^{\nu}$，则$dx_\mu$也是张量}
\item{如果定义$\partial^\mu\phi \equiv g^{\mu\nu}\partial_\mu\phi$，则$\partial^\mu\phi$是张量}
\item{如果$T^{\mu\nu}$是二阶张量，则$T^{\mu}_{\ \mu}$是标量}
\item{如果$A^{\mu}$和$B^{\mu}$都是矢量，则$A^\mu B^\nu$是二阶张量}
\end{itemize}

\ech
\end{frame}


\begin{frame}
\frametitle{\bch 张量指标的升降\ech}
\bch
通过前面的讨论，我们发现张量的坐标可以通过度规$g_{\mu\nu}$和其逆矩阵$g^{\mu\nu}$来升降
\ech
{\vskip 0.1in}

\begin{minipage}{0.3\textwidth}
\includegraphics[width=1in]{jugelizi.png}
\end{minipage}
\begin{minipage}{0.6\textwidth}
$$F_{\mu\nu} = g_{\mu\rho}g_{\nu\sigma}F^{\rho\sigma}$$
$$T^{\mu}_{\ \nu\rho} = g^{\mu\sigma} T_{\sigma\nu\rho}$$
$$g^{\mu}_{\ \nu} = g^{\mu\rho}g_{\rho\nu} = \delta^{\mu}_{\ \nu}$$
\end{minipage}
\end{frame}


\begin{frame}
\frametitle{\bch 张量的直积 \ech}
\bch
通过前面的讨论，我们发现一个m阶张量和一个n阶张量可以直接相乘生成m+n阶张量。
\ech
{\vskip 0.1in}

\begin{minipage}{0.3\textwidth}
\includegraphics[width=1in]{jugelizi.png}
\end{minipage}
\begin{minipage}{0.6\textwidth}
\bch
$dx^\mu dx^\nu$是二阶张量。

$g_{\mu\nu}g^{\rho\sigma}$ 是四阶张量
\ech
\end{minipage}
\end{frame}



\begin{frame}
\frametitle{\bch 张量指标的收缩\ech}
\bch
通过前面的讨论，我们发现张量的上下指标可以收缩，产生低两阶的张量。
\ech
{\vskip 0.1in}

\begin{minipage}{0.3\textwidth}
\includegraphics[width=1in]{jugelizi.png}
\end{minipage}
\begin{minipage}{0.6\textwidth}
$$dx^\mu dx_\mu = ds^2$$
$$g^{\mu}_{\ \mu} = 4$$
$$g^{\mu\rho}g_{\rho\nu} = \delta^{\mu}_{\ \nu}$$
\end{minipage}
\end{frame}



\begin{frame}
\frametitle{\bch课堂互动\ech}
\bch
\addfig{1}{think2.jpg}
如果$\phi$是标量，那么$\partial_\mu\partial_\nu\phi$是张量吗？
\ech
\end{frame}

\begin{frame}
  \chtitle{课后作业 (9月19日周二课后交）}
  \bch
  \bitem
\item[1]{自然单位制下力矩的量纲是什么（=质量的多少次方）？}
\item[2]{某散射过程的截面为$\sigma = 10^{-3}/m_W^2$, 其中$m_W \approx 80 \mathrm{GeV}$是$W^{\pm}$的粒子质量。试换算出以$\mathrm{cm}^2$为单位的截面值。}
\item[3]{设四维时空度规为$g_{\mu\nu}$，试化简$g_{\mu\nu}g^{\mu\nu}$。}
  \eitem
  \ech
\end{frame}

\end{document}
